\scanpage{65}\mspage{64}
Inversement, donnons nous, \deleted{$\unsure{F_i\in|F_i}$} $\unsure{(}\forall i\in I)$ un $F_i\in\mathcal F_i$, quand peut-on dire qu'il existe $F\in\mathcal F$ telle que $\forall i$, $F_i$ soit la $i$ composante de $\deleted{\unsure{\mathcal F}}\unsure{F_0|F}$ ?

Déjà, il le faudrait pour \unsure{le|les} cas où tous les $F_i$ sauf un nb fini sont \deleted{$\varnothing$} vides. On va supposer donc

\uline{$C4$} Si $x,y\in M$ appartiennent à \unsure{des|deux} comp. connexes disjointes de $M$, alors $x$ et $y$ sont compatibles.

\deleted{Cor} \inserted{C'est \unsure{néc|néc.} et suffisant pour la \deleted{\unsure{déd|val}} \unsure{validité|déduction} des}

\uline{Corollaire 1} L'image de $\mathcal F\to\prod\mathcal F_i$ contient les $(F_i)_{i\in I}$ tels que $F_i=\varnothing$ \deleted{sauf} pour \deleted{\unsure{pres}}que tout $i$.

\inserted{\uline{Cor. 2}} \deleted{\unsure{donc}} $\mathcal F\to\prod\mathcal F_i$ est une bijection dans \deleted{les} le \deleted{deux} cas où \deleted{suivants} \inserted{$I$ est} \deleted{\unsure{f}} fini.

\deleted{a) $I$ est fini}

\deleted{b) Toute figure est de t.f. \unsure{--} alors}

\uline{Cor 3} Supposons que les figures \deleted{\unsure{de}} $\in\mathcal F$ soient de t.f. \deleted{Alors il} \inserted{(resp. finies)}. Alors il en est de $\unsure{=|\text{même}}$ \deleted{\unsure{de}} pour $\mathcal F_i$, de plus $\mathcal F\unsure{\text{---}|\longrightarrow}\prod\mathcal F_i$ a comme image l'ens des $(F_i)$ tels que $F_i$ soit $\varnothing$ pour pres que tout $i$.
